% !PRINTED 162
\weberpage{162}
%<<CONT
gives, under this assumption, a rational expression for the common
root in terms of the coefficients.

If one substitutes this expression into $\psi(x)$, then, on clearing
away all denominators, one obtains an integral rational function
of the coefficients $a$, $b$, $c$, which we shall denote by $S$, and
\gesperrt{whose vanishing, in conjunction with $R = 0$, contains the necessary and
sufficient condition for the three equations (1) to have a common
root.}

Let us now replace the coefficients $a$, $b$, $c$ of the functions
(1) by integral rational functions of a second unknown
$y$, whose coefficients shall be denoted by $\alpha$, $\beta$, $\gamma$ and for the time being
regarded as independent variables.

$R$ and $S$ then become integral rational functions of
$\alpha$, $\beta$, $\gamma$ and $y$, and if we therefore regard the equations

\weq{2}{R = 0, \qquad S = 0}

as equations for $y$, then by the rules of the preceding
paragraph we can form their resultant $T$, which,
if it does not vanish identically, will be an integral rational
function of the coefficients $\alpha$, $\beta$, $\gamma$, and the equation

\weq{3}{T = 0}

is the necessary and sufficient condition for the
system of equations (1) to be satisfiable by one or more pairs of values $x$, $y$.
$T$ is called the \gesperrt{resultant of the system
(1).} % [source-original logical slip, confirmed 2026-08-25 against scan p. 187 / printed p. 162]
% The print says that eliminating x and y forms a final equation "for x", although x has just
% been eliminated. The facing translation preserves the printed variable rather than guessing.
One can therefore, from the three equations (1), by \gesperrt{elimination of the unknowns $x$ and $y$, form a final equation for
$x$.}

If one now replaces the $\alpha$, $\beta$, $\gamma$ in turn by integral rational
functions of a third unknown $z$, then $T$ passes over into an
equation for $z$. The functions $f$, $\varphi$, $\psi$ become functions
of the three unknowns $x$, $y$, $z$. The degree of equation (3) gives
the number of systems of values $x$, $y$, $z$ for which the three equations
(1) are simultaneously satisfied.

But if the coefficients are so constituted that $T$ vanishes
identically (for all values of $z$), then there are infinitely many
systems of values $x$, $y$, $z$ which satisfy the equations (1) simultaneously.

If one regards $x$, $y$, $z$ as Cartesian coordinates, then
the equations (1), each by itself, represent a surface.
The degree of the final equation $T = 0$ with respect to $z$ gives the
%CONT>>
